\chapter{Dispersion and Correlation}\label{ch:s2:dispersion-and-correlation}

An index option is a bet on the index's volatility, which depends on how volatile its members are and how much they move together. Cboe publishes the
correlation that option prices imply among the 50 largest S\&P 500 stocks: from 2006 to 2026 it averaged 37, reached 96.6 on 24 October 2008, and fell
to 2.93 in July 2024. Driessen, Maenhout and Vilkov found that correlation risk is priced: selling it, by selling index volatility against members'
volatility, earned a high alpha before frictions and none after them. On this chapter's synthetic market the same trade earns a Sharpe ratio of 0.65
made vega-neutral, and a pure short correlation swap 1.74, until a crash takes correlation from 0.12 to 0.83 in a month. The build is
\texttt{firm.dispersion}.

\section{Index volatility from member volatilities}

An index's variance is the weighted sum of its members' variances and covariances: $\sigma_I^2 = \sum_i w_i^2\sigma_i^2 + \sum_{i \ne j} w_i w_j
\rho_{ij}\sigma_i\sigma_j$. Replacing every pairwise correlation by one average $\bar\rho$ and solving gives the average correlation that makes the
members' volatilities consistent with the index's:
\[
\bar\rho = \frac{\sigma_I^2 - \sum_i w_i^2\sigma_i^2}{\sum_{i \ne j} w_i w_j \sigma_i\sigma_j}.
\]
Cboe's implied correlation indices apply this decomposition to the implied volatilities of the S\&P 500 and of its 50 largest members.

\begin{definition}[Realised correlation]\label{def:s2:dispersion-and-correlation:realised}
\emph{Realised correlation}\index{realised correlation} is the average correlation among an index's members computed from realised variances over a
period with the same decomposition as implied correlation: the index's realised variance against the members' realised variances.
\end{definition}

\texttt{firm.dispersion} (\cref{lst:s2:dispersion-and-correlation:legs}) computes both. On the synthetic market the index is the equal-weighted
average of thirty members; members' implied variance is their expected variance times 1.10 and the index's times 1.30, so the index's options are dearer
than its members' and implied correlation sits above what follows: 0.311 on average against a realised 0.252, higher in 77\% of months.

\section{The correlation risk premium}

\begin{definition}[Correlation risk premium]\label{def:s2:dispersion-and-correlation:premium}
The \emph{correlation risk premium}\index{correlation risk premium} is the difference between the average correlation implied by index and member
option prices and the correlation then realised; it is positive on average because correlation rises in crashes, when diversification is most needed.
\end{definition}

Driessen, Maenhout and Vilkov used S\&P 100 index options and options on all its components: correlation risk was priced, it explained the
cross-section of index and single-stock option returns, and a strategy selling it had a high alpha for an investor without frictions. With realistic
trading costs the premium could not be exploited, which they read as a limit to arbitrage. Cboe's own comparison with a smoothed \omterm{def:s2:dispersion-and-correlation:realised}{realised correlation}
of the top 50 stocks finds realised usually slightly below implied.

The real index (\cref{fig:s2:dispersion-and-correlation:cor}) shows why selling correlation is uncomfortable: COR1M averaged 37.0 from January 2006 to
September 2026, with 5\% and 95\% quantiles of 10.8 and 69.7, and averaged 70.4 in October 2008 and 72.5 in March 2020. On 22 September 2026 it stood at
7.59; its low was 2.93, on 12 July 2024. Cboe's S\&P 500 dispersion index, published since
June 2014, averaged 25.5 and peaked at 58.9 on 18 March 2020.

\begin{omfigure}
\begin{tikzpicture}
\begin{axis}[width=11cm,height=5.2cm,xlabel={\small year},ylabel={\small COR1M, monthly mean},tick label style={font=\footnotesize},
  xmin=2006,xmax=2027,ymin=0,ymax=100,grid=major,grid style={gray!25},table/col sep=comma,/pgf/number format/1000 sep={}]
\addplot[omThm,thick] table[x=year,y=cor1m]{figdata/strategies-2/02-dispersion-and-correlation/cor1m.csv};
\end{axis}
\end{tikzpicture}

\omcaption{Cboe's 1-Month Implied Correlation Index (the average correlation among the S\&P 500's 50 largest members implied by index and single-stock
options), monthly means, January 2006 to September 2026. Derived from Cboe's index history; the raw series is not redistributed. Data:
\texttt{s2\_fetch\_cor}.}\label{fig:s2:dispersion-and-correlation:cor}
\end{omfigure}

\section{Dispersion trades and their weights}

A dispersion trade (Book~5, chapter~17) sells index volatility and buys members' volatility. The chapter's books use one-month variance swaps, rebuilt
monthly (\cref{lst:s2:dispersion-and-correlation:books}): short one unit of index variance, long members' variance on notionals set one of two ways.

\begin{definition}[Dispersion weighting]\label{def:s2:dispersion-and-correlation:weighting}
\emph{Dispersion weighting}\index{dispersion weighting} is the choice of the members' notionals in a dispersion trade: equal to their index weights, so
that the trade is short correlation and long the members' specific variance; or scaled so that the members' total vega equals the index's
(vega-neutral), so that a common rise in volatility leaves the book flat and its P\&L comes mainly from correlation.
\end{definition}

\begin{center}
\small
\begin{tabular}{@{}lcccc@{}}
\toprule
239 months & Sharpe ratio & skewness & worst month (sd) & months up\\
\midrule
members on index weights & $-0.02$ & 2.67 & $-1.6$ & 36\%\\
vega-neutral & 0.65 & $-5.87$ & $-11.3$ & 70\%\\
short correlation swap & 1.74 & $-1.67$ & $-6.0$ & 77\%\\
\bottomrule
\end{tabular}
\end{center}

On index weights the book is long more member variance than it is short index variance, and it pays the members' own premium: it loses a little on
average and gains in crashes (positive skewness). Made vega-neutral it earns the correlation premium with the familiar short-volatility shape. The
correlation swap, which pays implied minus \omterm{def:s2:dispersion-and-correlation:realised}{realised correlation} directly, isolates the premium best; its Sharpe ratio is the highest and its worst month
is still six standard deviations.

\section{When correlation goes to one}

The three planted crashes are the tests (\cref{fig:s2:dispersion-and-correlation:books}). In the first, implied correlation was 0.12 when the month began
and \omterm{def:s2:dispersion-and-correlation:realised}{realised correlation} 0.83: the vega-neutral book lost 11.3 of its monthly standard deviations and the correlation swap 0.71 of a unit. In the second,
0.12 against 0.57 (3.0 standard deviations and 0.45). In the third, which began when implied correlation was already 0.73, the loss was under one
standard deviation. A crash is a correlation event: every member falls together, the index falls as far as its members, and the diversification the
index option was priced to lack arrives.

\begin{omfigure}
\begin{tikzpicture}
\begin{axis}[width=11cm,height=5.2cm,xlabel={\small year},ylabel={\small cum.\ P\&L (sd units)},tick label style={font=\footnotesize},
  xmin=0,xmax=20,ymin=-20,ymax=130,grid=major,grid style={gray!25},legend pos=north west,legend style={font=\scriptsize},table/col sep=comma]
\addplot[omDef,thick] table[x=year,y=vega]{figdata/strategies-2/02-dispersion-and-correlation/books.csv};
\addplot[omThm,thick] table[x=year,y=corrswap]{figdata/strategies-2/02-dispersion-and-correlation/books.csv};
\legend{vega-neutral dispersion,short correlation swap}
\end{axis}
\end{tikzpicture}

\omcaption{The synthetic vega-neutral dispersion book and short correlation swap, each cumulated in units of its monthly standard deviation, over
twenty years with crashes at years 6.0, 12.7 and 17.5. Data: \texttt{s2\_dispersion.monthly}.}\label{fig:s2:dispersion-and-correlation:books}
\end{omfigure}

\section{Strategy files}

\begin{strategyfile}{Vega-weighted dispersion}\label{strat:s2:dispersion-and-correlation:vega}
\sfield{Who pays you, and why} Buyers of index protection, who pay for the correlation that crashes bring.
\sfield{Instruments and venues} Index options or variance swaps against single-stock options or variance swaps on the largest members.
\sfield{Signal} Implied against forecast correlation.
\sfield{Sizing and execution} Members' vega matched to the index's; monthly or quarterly rolls; members limited to liquid names.
\sfield{Costs} Single-stock option spreads, the main friction.
\sfield{How it dies} Crashes; frictions that absorb the premium.
\sfield{Horizon, capacity, infrastructure} Months; single-stock surfaces for dozens of names.
\sfield{Backtest honestly} Single-stock option costs; crash months in the sample.
\sfield{Sources} Driessen, Maenhout and Vilkov (2009); this chapter: 0.65 and an 11-standard-deviation crash month.
\end{strategyfile}

\begin{strategyfile}{Theta-weighted dispersion}\label{strat:s2:dispersion-and-correlation:theta}
\sfield{Who pays you, and why} As for vega-weighted dispersion.
\sfield{Instruments and venues} Index and single-stock straddles.
\sfield{Signal} As above.
\sfield{Sizing and execution} Members' theta matched to the index's, so that the book's time decay is neutral and its P\&L is gamma against gamma.
\sfield{Costs} Delta hedging of many options.
\sfield{How it dies} As above.
\sfield{Horizon, capacity, infrastructure} Weeks to months; a hedging engine.
\sfield{Backtest honestly} Hedging costs for every member.
\sfield{Sources} No performance figure verified.
\end{strategyfile}

\begin{strategyfile}{Correlation swap short}\label{strat:s2:dispersion-and-correlation:corrswap}
\sfield{Who pays you, and why} Buyers of correlation protection, often structured-product desks (chapter~28).
\sfield{Instruments and venues} Over-the-counter correlation swaps.
\sfield{Signal} Strike against forecast correlation.
\sfield{Sizing and execution} Sized by the crash-month loss.
\sfield{Costs} Dealer spreads; counterparty exposure.
\sfield{How it dies} Crashes.
\sfield{Horizon, capacity, infrastructure} Months to a year; an over-the-counter relationship.
\sfield{Backtest honestly} \omterm{def:s2:dispersion-and-correlation:realised}{Realised correlation} measured as the contract defines it.
\sfield{Sources} This chapter: 1.74 before costs, a loss of 0.71 of a unit in the first crash.
\end{strategyfile}

\begin{strategyfile}{Sector dispersion}\label{strat:s2:dispersion-and-correlation:sector}
\sfield{Who pays you, and why} As for index dispersion, within a sector.
\sfield{Instruments and venues} Sector ETF options against their largest members' options.
\sfield{Signal} Sector implied correlation against its history.
\sfield{Sizing and execution} Vega-neutral within sectors.
\sfield{Costs} Wider spreads on sector options.
\sfield{How it dies} Sector-wide shocks.
\sfield{Horizon, capacity, infrastructure} Months; small capacity.
\sfield{Backtest honestly} Liquidity of sector options.
\sfield{Sources} Cboe's sector implied correlations (as a data source); no performance figure verified.
\end{strategyfile}

\section{Tutorial: when everything moves together}\label{tut:s2:dispersion-and-correlation:tut}

\begin{tutorial}
\textbf{Goal.} Measure implied and \omterm{def:s2:dispersion-and-correlation:realised}{realised correlation} on the synthetic index and its members, run three dispersion books through the planted
crashes, and read Cboe's correlation and dispersion indices. \textbf{End state:} the table and the two figures.
\begin{tutsteps}
\item \textbf{Correlation and the legs}.
\omcode{code/firm/dispersion/firm_dispersion.py}{21}{47}{Average correlation, realised variance and the dispersion P\&L.}\label{lst:s2:dispersion-and-correlation:legs}
\item \textbf{The monthly books}.
\omcode{code/strategies-2/02-dispersion-and-correlation/python/s2_dispersion.py}{35}{58}{Implied and realised correlation and three books, month by month.}\label{lst:s2:dispersion-and-correlation:books}
\item \textbf{Run} \texttt{s2\_fetch\_cor.py} once, then \texttt{monthly()}, \texttt{crashes()} and \texttt{fig\_dispersion.py}.
\end{tutsteps}
\textbf{What to change next.} Trade only the twenty largest members and measure the tracking of the index; add a correlation skew so that crash
correlation is priced; hedge the vega-neutral book with a small long correlation position before calm periods end.
\end{tutorial}

\section{Build: dispersion}\label{bld:s2:dispersion-and-correlation:dispersion}

\begin{build}
\bfield{Purpose} Implied and realised average correlation, dispersion books on variance swaps and correlation swaps.
\bfield{Interface} \texttt{average\_corr(var\_index, var\_members, w)}, \texttt{realised\_var(R, start, n)}, \texttt{dispersion\_pnl(iv\_index,
iv\_members, rv\_index, rv\_members, w, vega\_neutral)}, \texttt{corr\_swap\_pnl(implied, realised)}.
\bfield{Rules} The same decomposition for implied and \omterm{def:s2:dispersion-and-correlation:realised}{realised correlation}; variance swap legs in annual variance units.
\bfield{Acceptance tests} \texttt{code/firm/dispersion/tests/}: a common correlation recovered exactly; zero P\&L when every leg realises its implied
variance; the vega-neutral notionals by hand.
\bfield{Stretch} Straddle legs with hedging; correlation skew; sector baskets.
\end{build}

\begin{omsources}
\item J.~Driessen, P.~J.~Maenhout and G.~Vilkov, ``The price of correlation risk: evidence from equity options'', \emph{Journal of Finance} 64(3), 2009.
\item Cboe, S\&P 500 Implied Correlation Index methodology and implied correlation notes.
\item Cboe Global Markets, daily histories of COR1M and DSPX.
\end{omsources}

\section{Exercises}

\begin{exercise}[$\star$]\label{exo:s2:dispersion-and-correlation:1}
An index of many equal-weighted members has implied volatility 20\% and each member 30\%. What average correlation does that imply?
\end{exercise}

\begin{exercise}[$\star$]\label{exo:s2:dispersion-and-correlation:2}
For a vega-neutral dispersion trade with index implied volatility 20\% and members at 30\%, by what factor are the members' notionals scaled from
their index weights?
\end{exercise}

\begin{exercise}[$\star$]\label{exo:s2:dispersion-and-correlation:3}
Why does the book on index weights have positive skewness?
\end{exercise}

\begin{exercise}[$\star\star$]\label{exo:s2:dispersion-and-correlation:4}
A short correlation swap was struck at 0.12 and correlation realised 0.83. What does it lose per unit, and why did the strike look reasonable when set?
\end{exercise}

\begin{exercise}[$\star\star$]\label{exo:s2:dispersion-and-correlation:5}
Why could Driessen and co-authors find a priced correlation premium that could not be exploited?
\end{exercise}

\begin{exercise}[$\star\star$]\label{exo:s2:dispersion-and-correlation:6}
COR1M fell to 2.93 in July 2024. What does a low implied correlation say about index options against single-stock options, and what risk does
selling correlation there carry?
\end{exercise}

\begin{exercise}[$\star\star\star$]\label{exo:s2:dispersion-and-correlation:7}
\emph{Coding.} Run \texttt{monthly(0.30)}, which sets the members' premium equal to the index's. What happens to the correlation premium and
the books?
\end{exercise}

\begin{exercise}[$\star\star\star$]\label{exo:s2:dispersion-and-correlation:8}
\emph{Find the flaw.} ``Implied correlation has exceeded realised in 77\% of months, so a short correlation swap is a high-probability trade we can
lever.''
\end{exercise}

\section{Problem: When Everything Moves Together}

\begin{problem}[{Weekend problem --- selling correlation}]\label{pb:s2:dispersion-and-correlation:1}
The chapter's synthetic index, Cboe's indices and the public record.

\medskip\noindent\textbf{Part I --- Correlation.}
\begin{enumerate}
\item Write the average-correlation formula and explain it.
\item Define \omterm{def:s2:dispersion-and-correlation:realised}{realised correlation} and the \omterm{def:s2:dispersion-and-correlation:premium}{correlation risk premium}.
\item How does Cboe compute its implied correlation index?
\item What did Driessen, Maenhout and Vilkov find?
\end{enumerate}

\medskip\noindent\textbf{Part II --- The real indices.}
\begin{enumerate}[resume]
\item Give COR1M's average, range and its 2008 and 2020 levels.
\item What does DSPX measure, and what were its statistics?
\item Why is a record low in implied correlation not a free trade?
\item What did Cboe find comparing implied and \omterm{def:s2:dispersion-and-correlation:realised}{realised correlation}?
\end{enumerate}

\medskip\noindent\textbf{Part III --- The books.}
\begin{enumerate}[resume]
\item Define \omterm{def:s2:dispersion-and-correlation:weighting}{dispersion weighting} and describe the two weightings.
\item Give the three books' Sharpe ratios and skewness.
\item Why does the book on index weights lose on average?
\item Why is the correlation swap the cleanest exposure?
\end{enumerate}

\medskip\noindent\textbf{Part IV --- The verdict.}
\begin{enumerate}[resume]
\item State the \emph{named result}: the dispersion book's return from the correlation premium and its loss when correlation jumps.
\item What happened in each of the three crashes?
\item Why did the third crash cost less?
\item How would you size a dispersion book?
\item How would you backtest one honestly?
\item Which strategy file is most exposed to single-stock option costs?
\item How does this chapter relate to chapter~1?
\item In one sentence: what does a dispersion trader sell?
\end{enumerate}
\end{problem}

\section{Interview questions}

\begin{interviewq}[$\star$ \iqroles{researcher}]\label{iq:s2:dispersion-and-correlation:1}
What is implied correlation, and how is it computed?
\end{interviewq}

\begin{interviewq}[$\star\star$ \iqroles{trader}]\label{iq:s2:dispersion-and-correlation:2}
Walk through a dispersion trade and the risks it carries.
\end{interviewq}

\begin{interviewq}[$\star\star$ \iqroles{researcher}]\label{iq:s2:dispersion-and-correlation:3}
Why might index options be expensive relative to single-stock options?
\end{interviewq}

\begin{interviewq}[$\star\star$ \iqroles{risk}]\label{iq:s2:dispersion-and-correlation:4}
How would you stress a dispersion book?
\end{interviewq}

\begin{interviewq}[$\star\star$ \iqroles{developer}]\label{iq:s2:dispersion-and-correlation:5}
What data does a dispersion desk need every day, and what can go wrong with it?
\end{interviewq}

\begin{interviewq}[$\star\star\star$ \iqroles{researcher}]\label{iq:s2:dispersion-and-correlation:6}
For $N$ equal-weighted members with equal volatility $\sigma$ and average correlation $\rho$, show that the index variance is $\sigma^2(\rho + (1 -
\rho)/N)$, and derive the implied correlation from index and member volatilities for large $N$.
\end{interviewq}
